% ===================================================================
%  اللوحة نمرة ١٣ — اللوكاندة. --- المطعم. القهوة.
%
%  Traduction, faite en 2026, en arabe egyptien ecrit, sur l'ido de
%  texto/io. Regles de la colonne : voir 10-lawha-01.tex. Une seule
%  numerotation. Deux notes d'auteur, toutes deux en gras-italique.
%
%  L'ASCENSEUR (9) EST الأسانسير, ET C'EST LE MEME FAIT QUE 𨋢 EN
%  CANTONAIS. Les deux colonnes ont trouve, au meme endroit du meme
%  tableau, que le mot de l'ascenseur est un emprunt europeen entre
%  par le port : lift a Hong Kong, ascenseur a Alexandrie. L'arabe
%  standard dit مصعد, qui est une traduction.
%
%  ET LE POURBOIRE (59) EST البقشيش, QUE L'ANGLAIS A EMPRUNTE A SON
%  TOUR sous la forme « baksheesh ». C'est le seul mot de toute la
%  colonne qui ait fait le voyage dans l'autre sens.
%
%  LES JEUX DE LA SALLE DE CAFE SONT TOUS EGYPTIENS, ET DEUX D'ENTRE
%  EUX TOMBENT EXACTEMENT :
%
%    طاولة (81)    le trictrac — c'est le nom du jeu en Egypte, et
%                    c'est bien le meme jeu
%    كوتشينة (82)  les cartes  de l'italien
%    ضامة (79)     les dames
%    دومينو (80)   les dominos
%    بلياردو (75)  le billard
%
%  LE RESTE DU SEJOUR :
%
%    كاسّة (15)     la caisse    de l'italien cassa
%    بوّاب (18)     le concierge
%    جواب (22)     la lettre    MSA : رسالة
%    كرباج (93)    le fouet
%    تلغرافجي (35) le porteur de depeches — le septieme -جي
%    جمبري         les crevettes
%    كرشة          les tripes
%
%  ET LES TROIS REPAS SE SEPARENT COMME AU QUEBEC ET A CANTON, MAIS
%  AVEC LEURS PROPRES MOTS : الفطار le matin, الغدا a midi, العشا le
%  soir. L'ido distingue lui-meme « dejuneto » et « dejuno » dans ce
%  tableau — c'est ce qui avait fait la preuve pour fr-CA — et
%  l'egyptien les rend par الفطار et الغدا sans avoir a choisir.
%
%  L'ANNEXION ARABE PASSE TOUS LES PIEGES DU TABLEAU SANS UNE
%  REPRISE : la chambre et son etage, le cocher de l'omnibus, la
%  malle du voyageur, les nuages de l'horizon, l'enseigne du hall.
%  Les cinq avaient demande une reecriture en cantonais.
% ===================================================================

%%K t13-noto-1 noto
\VUnotes{143.2mm}{%
\textit{\VUgras{(*) بالنسبة لتفاصيل الأوضة، بُصّ اللوحة نمرة ٦.}}%
}

%%K t13-apar-1 apar
\VUsaut{3.0mm}
\VUtitre{15.0pt}{60}{اللوحة نمرة ١٣}
\VUsaut{1.6mm}
\VUpk{10.2pt}{40}{اللوكاندة. --- المطعم.}
\VUsaut{2.0mm}
\VUpk{10.2pt}{40}{القهوة.}
\VUsaut{3.4mm}

%%K t13-01-1 p
1. --- وصلت إمبارح بالليل رويان، ونزلت في «\textit{لوكاندة المحيط}» اللي كان
صاحب لي نصحني بيها.

%%K t13-01-2 p
وإدّوني \VUgras{أوضة} \textsuperscript{(2)} حلوة، في \VUgras{الدور التاني}
\textsuperscript{(1)}، ونمت فيها كويس أوي (*).

%%K t13-02-1 p
2. --- والنهارده الصبح، وأنا خارج من أوضتي، اللي \VUgras{الشغّالة}
\textsuperscript{(3)} كانت جابت لي فيها الفطار، دخلت \VUgras{أوضة التدخين}
\textsuperscript{(5)} اللي بتتستعمل كمان \VUgras{أوضة قراية} \textsuperscript{(4)}،
علشان أقرا \VUgras{الجرايد} \textsuperscript{(6)} وأنا بشرب \VUgras{سيجارة}
\textsuperscript{(8)}. وبعد ما قريت، نزلت بـ\VUgras{الأسانسير} \textsuperscript{(9)}
ورحت أتمشّى في المدينة.

%%K t13-03-1 p
3. --- وعلشان نسيت أسأل \VUgras{بتاع الحسابات} \textsuperscript{(12)} بتاع اللوكاندة
الأكل بيتقدّم على \VUgras{الترابيزة العمومية} \textsuperscript{(11)} الساعة كام، رجعت
بدري واستغلّيت ده علشان أستحمّى في \VUgras{حمّام} \textsuperscript{(10)} اللوكاندة.
وبعدين رحت تاني لـ\VUgras{الكاسّة} \textsuperscript{(15)} علشان أكلّم واحد من أصحابي
في التليفون. بس \VUgras{التليفون} \textsuperscript{(13)} كان مشغول؛ ولما
\VUgras{الست بتاعة الكاسّة} \textsuperscript{(14)} شافت حيرتي --- وهي كانت بتسجّل
\VUgras{فاتورة} \textsuperscript{(17)} في \VUgras{دفتر النزلاء} \textsuperscript{(16)}
--- رجت \VUgras{البوّاب} \textsuperscript{(18)} إنه يبعت \VUgras{الصبي}
\textsuperscript{(19)} يعمل لي المشوار \VUgras{بالعجلة} \textsuperscript{(20)}.

%%K t13-04-1 p
4. --- وشكرتها وخرجت علشان أدّي بقشيش لـ\VUgras{الشيّال} \textsuperscript{(24)} (اللي
بيشيل التقيل) اللي

%%K t13-04-1 p suite
جاب لي من المحطة على \VUgras{عربية إيده} \textsuperscript{(25)} \VUgras{علبة
البرنيطة} \textsuperscript{(26)}، و\VUgras{شنطتي} \textsuperscript{(27)}،
و\VUgras{شنطة سفري} \textsuperscript{(28)}. و\VUgras{البوسطجي} \textsuperscript{(21)}
اللي كان لسه واصل، إدّاني \VUgras{جواب} \textsuperscript{(22)}.

%%K t13-05-1 p
5. --- وفضلت شوية كمان على عتبة الباب، جنب \VUgras{صندوق البوستة}
\textsuperscript{(23)}، أتفرّج على الحركة في الشارع. و\VUgras{عربجي}
\textsuperscript{(30)} \VUgras{الأومنيبوس} \textsuperscript{(29)} كان بيشحن على عربيته
\VUgras{شنطة} \textsuperscript{(31)} \VUgras{نزيل} \textsuperscript{(32)} اللي
\VUgras{صاحب اللوكاندة} \textsuperscript{(33)} كان بيودّعه. و\VUgras{تلغرافجي}
\textsuperscript{(35)} صغيّر كان بيوصّل على مهله \VUgras{تلغراف} \textsuperscript{(34)}
لواحد من زباين اللوكاندة؛ و\VUgras{سايح} \textsuperscript{(36)}، و\VUgras{كاميرته}
\textsuperscript{(38)} على كتفه، كان بيبصّ في \VUgras{دليله} \textsuperscript{(37)}.

%%K t13-tit-1 sub
\VUsaut{4.0mm}
\VUpk{10.2pt}{40}{وقت الغدا.}
\VUsaut{3.0mm}

%%K t13-06-1 p
6. --- وقدام السور، \VUgras{رسول} \textsuperscript{(39)} مش مهموم كان بيغفّى بين
\VUgras{شيّالته} \textsuperscript{(40)} و\VUgras{صندوق البويا} \textsuperscript{(41)}
بتاعه، ولمّا \VUgras{غسّالة} \textsuperscript{(42)} صغيّرة كانت شايلة \VUgras{مقطف}
\textsuperscript{(43)} غسيل ضحكت على الوش الجادّ بتاع \VUgras{رئيس الصالة}
\textsuperscript{(44)} اللي كان واقف جنب الباب.

%%K t13-07-1 p
7. --- ولما شفت \VUgras{الشيف} \textsuperscript{(45)} خارج من \VUgras{مطبخه}
\textsuperscript{(47)} اللي في \VUgras{البدروم} \textsuperscript{(48)}، علشان يبعت
\VUgras{غاسل الأطباق} \textsuperscript{(46)} يعمل مشوار، افتكرت إن ساعة الغدا قرّبت
ودخلت المطعم، وأنا بعدّي الحوش اللي فيه \VUgras{سوّاق} \textsuperscript{(50)} بيركن
\VUgras{عربيته} \textsuperscript{(49)}، ولمّا \VUgras{ميكانيكي} \textsuperscript{(52)}
بيصلّح، على حسب كلام \VUgras{بتاع العجلة} \textsuperscript{(53)}، \VUgras{تريسيكل
بنزين} \textsuperscript{(51)} كان لسه اتخبط.

%%K t13-08-1 p
8. --- وسألت \VUgras{صبي} \textsuperscript{(54)} صغيّر كان شايل \VUgras{كبايات بيرة}
\textsuperscript{(55)}، الترابيزة العمومية تحت الفيرندة ولا لأ،

%%K t13-noto-2 noto
\VUnotes{143.2mm}{%
\textit{\VUgras{(*) بمساعدة لستة الأصناف اللي في القاموس، المدرِّس هيقدر بسهولة
يغيّر في المنيو اللي تحت.}}%
}

%%K t13-08-1 p suite
ولما ردّ بالإيجاب، أخدت مكاني على واحدة من الترابيزات وطلبت أكلة ممتازة.

%%K t13-tit-2 sub
\VUsaut{4.0mm}
\VUpk{10.2pt}{40}{غدا ممتاز.}
\VUsaut{3.0mm}

%%K t13-09-1 p
9. --- والجرسون جاب لي \VUgras{منيو} (*) الغدا ولستة النبيت. واخترت وطلبت
\VUgras{جمبري} مع \VUgras{زبدة} \VUgras{كمقبّلات}، و\VUgras{كرشة على طريقة كاين}
كمقدّمة. وما أكلتش \VUgras{سمك}، لكن طلبت جزء من \VUgras{فخدة} خروف صغير،
و\VUgras{صنف خضار}: \VUgras{سبانخ} بالقشدة. وبعدين، علشان مافيش صنف من
\VUgras{التحلاية} عجبني، طلبت حاجات بعد الأكل: \VUgras{جبنة}، و\VUgras{فاكهة}،
و\VUgras{بسكوت}. وضفت كمان \VUgras{نبيت} سان إميليون ممتاز، وأنا مبسوط أوي من
أكلتي سبت الترابيزة بعد ما إدّيت \VUgras{الجرسون} \textsuperscript{(57)}
\VUgras{بقشيش} \textsuperscript{(59)} كويس.

%%K t13-10-1 p
10. --- ولما شافني مستعدّ أمشي، \VUgras{المدير} \textsuperscript{(60)} اقترح عليّ
أطلع البلكونة علشان أتفرّج على منظر \VUgras{المحيط} \textsuperscript{(62)} الرائع.
ومن بعيد بانت \VUgras{قوارب صيد} \textsuperscript{(63)}، و\VUgras{مراكب شراعية}
\textsuperscript{(64)}، و\VUgras{باخرة} \textsuperscript{(65)} ضخمة، اللي شكلها الأسود
باين على \VUgras{السحاب} \textsuperscript{(67)} الأبيض بتاع \VUgras{الأفق}
\textsuperscript{(66)}. ونسمة خفيفة خلّت \VUgras{البيرق} \textsuperscript{(70)} المركّب
فوق \VUgras{يافطة} \textsuperscript{(69)} \VUgras{الصالة} \textsuperscript{(68)} يتموّج.

%%K t13-tit-3 sub
\VUsaut{3.0mm}
\VUpk{10.2pt}{40}{التسلية.}
\VUsaut{3.0mm}

%%K t13-11-1 p
11. --- والمنظر كان مشوّق أوي، لدرجة إني قرّرت أطلع بكرة على تراس القهوة وأنا
شايل \VUgras{نضّارة معظّمة} \textsuperscript{(71)} أو \VUgras{نضّارة مكبّرة}
\textsuperscript{(72)}.

%%K t13-12-1 p
12. --- ولما سبت البلكونة، رحت قهوة اللوكاندة علشان أشرب \VUgras{حاجة}
\textsuperscript{(73)} وألعب دور \VUgras{بلياردو} \textsuperscript{(75)}. وعدّيت من غير
ما أقف صالة القهوة اللي فيها زباين كتير بيلعبوا \VUgras{شطرنج} \textsuperscript{(78)}،
أو \VUgras{ضامة} \textsuperscript{(79)}، أو \VUgras{دومينو} \textsuperscript{(80)}، أو
\VUgras{طاولة} \textsuperscript{(81)}، أو \VUgras{كوتشينة} \textsuperscript{(82)}، وطلعت
على طول \VUgras{أوضة البلياردو} \textsuperscript{(74)}.

%%K t13-13-1 p
13. --- ولما الدور خلص، نزلت الدور الأرضي وطلبت من الجرسون \VUgras{ظرف}
\textsuperscript{(84)}، و\VUgras{ورق} \textsuperscript{(83)}، و\VUgras{طابع}
\textsuperscript{(85)}، و\VUgras{قلم} \textsuperscript{(87)}، و\VUgras{حبر}
\textsuperscript{(86)} علشان أكتب جواب.

%%K t13-13-2 p
وبعدين ناديت \VUgras{عربجي} \textsuperscript{(92)} كانت \VUgras{عربيته}
\textsuperscript{(88)} واقفة قدام القهوة. وسألته عن السعر بالساعة أو بالمشوار، ولما
اتّفقنا على السعر، فتح لي \VUgras{باب العربية} \textsuperscript{(89)} وركّبني. وطلع
تاني على \VUgras{الدكّة} \textsuperscript{(91)}، وحرّك الخيل بسرعة بـ\VUgras{الكرباج}
\textsuperscript{(93)}، وقمنا في نزهة لذيذة.

\VUsaut{6.0mm}
\VUfilet{22mm}
